\scanpage{338}\mspage{254}
% l. 1–4
diagramme
\diagram{\begin{tikzcd}[ampersand replacement=\&,column sep=normal]
 F_{\alpha}(j)\arrow[r,"u_{\alpha}(j)"]\arrow[d,"F_{\alpha}(\lambda)"',no head]\arrow[dr,phantom,"!" description]
 \& F(j)\arrow[r,"v(j)",no head]\arrow[d,"F(\lambda)"']
 \& X(j)\arrow[r,"w(j)"]\arrow[d,phantom,"\text{comm.}" description]
 \& G(j)\arrow[d,"G(\lambda)"]\\
 F_{\alpha}(j')\arrow[r,"u_{\alpha}(j')",no head]\& F(j')\arrow[r,"v(j')",no head]\&
 X(j')\arrow[r,"w(j')",no head]\&G(j')\,.
\end{tikzcd}}

% l. 5–8
Le lemme 4.10.5. étant prouvé, \unsure{tout|il} revient donc à prouver la condition $\mathrm F1$, et son complément $\mathrm F'1$, qui se \unsure{prouve|prouvent} dans la $\overline{\mathrm m}$ haleine.

% l. 9–10
4.10.6. On se donne donc le système d'objets
\[
 X(j)\in\operatorname{Ob}\unsure{\mathcal E|\mathscr E}_j\qquad (j\in\operatorname{Ob}J), \tag*{(4.10.6.1)}
\]
% l. 11–13
et le cas échéant, un objet $F$ de $\underline{\Gamma}(\mathcal E/J)$ et un système de flèches
\[
 F(j)\overset{u(j)}{\unsure{\mathbin{-}|\longrightarrow}}X(j), \tag*{(4.10.6.2)}
\]
% l. 14–16
et on se propose de construire $G$ dans $\underline{\Gamma}(\mathcal E/J)$, et un système de flèches
\[
 X(\unsure{j|j'})\xrightarrow{v(j)}G(j), \tag*{(4.\unsure{10.6|10.26}.3)}
\]
% l. 17–19
de telle façon que (si on a la donnée supplémentaire des $u(j)$) les flèches composées $\deleted{\unsure{X}}F(j)\overset{u(j)}{\unsure{\mathbin{-}|\longrightarrow}}X(j)\overset{v(j)}{\unsure{\mathbin{-}|\longrightarrow}}G(j)$
% archive sheet number: 337
